\documentclass[11pt,a4paper]{article}
\usepackage[margin=1in]{geometry}
\usepackage{amsmath,amssymb,amsfonts,amsthm}
\usepackage{bm}
\usepackage{booktabs}
\usepackage{enumitem}
\usepackage{hyperref}
\usepackage{cleveref}

% ── Notation macros (consistent with report 1 and Pythonic-DISORT.ipynb) ───
\newcommand{\R}{\mathbb{R}}
\newcommand{\tr}{\operatorname{tr}}
\newcommand{\diag}{\operatorname{diag}}
\newcommand{\cond}{\operatorname{cond}}
\newcommand{\expm}{\operatorname{expm}}
\newcommand{\vect}[1]{\bm{#1}}
\newcommand{\mat}[1]{\mathbf{#1}}
\newcommand{\tbot}{\tau_{\mathrm{bot}}}
\newcommand{\tsub}{\tau_{\mathrm{sub}}}
\newcommand{\dd}{\,\mathrm{d}}
\newcommand{\kmax}{k_{\max}}
\newcommand{\kmin}{k_{\min}}
\newcommand{\Rinf}{\mat{R}_\infty}

\theoremstyle{definition}
\newtheorem{proposition}{Proposition}
\newtheorem{remark}{Remark}
\newtheorem{definition}{Definition}

\title{Diffusion-Domain Decomposition for the Star-Product Magnus Integrator:\\
Eigenvalue Structure, Matrix Riccati Equation, and Implicit Solvers}
\author{Technical Report~II --- Extension to the Pythonic-DISORT\\
Comprehensive Documentation}
\date{March 2026}

\begin{document}
\maketitle

\begin{abstract}
The star-product Magnus integrator of Report~I~\cite{report1} is
unconditionally stable for arbitrary optical depth, but the step
count $K$ is governed by the constraint $h\cdot\lambda_{\max}\lesssim 1$,
where $\lambda_{\max}\sim 1/\mu_{\min}$ is the fastest (ballistic)
eigenvalue of the coefficient matrix.  In the diffusion regime---thick,
highly scattering atmospheres where the physics is governed by the much
smaller diffusion eigenvalue $\kmin$---this constraint forces
$K\propto\lambda_{\max}\tbot$, even though $\kmin\ll\lambda_{\max}$.
This report investigates whether the resulting eigenvalue gap $\gamma
=\kmax/\kmin$ can be exploited via domain decomposition combined with
stiff-aware solvers.

We develop the matrix Riccati equation (invariant imbedding) formulation
for the reflection operator, prove that it is stiff with the same
eigenvalue gap $\gamma$, and show that explicit integrators (RK4) offer
\emph{no speedup} over the Magnus approach.  Implicit solvers
(Radau~IIA) handle the stiffness efficiently: for $N\!=\!8$ streams
($\text{NQuad}=16$), the Radau solver requires $\approx 5000$ function
evaluations independent of $\omega$, compared to $K\approx 914$ Magnus
steps.  The wall-time advantage is modest (${\sim}1.5\times$) and does
not scale with $\gamma$, because the Riccati Jacobian is
$N^2\!\times\!N^2$ and its formation and factorisation dominate the cost.
A three-domain star-product decomposition (non-diffusive/diffusive/non-diffusive)
is validated as exact to discretisation accuracy.

A Rosenbrock method that exploits the Sylvester structure of the
Riccati Jacobian reduces each implicit step to an $N\!\times\!N$
Sylvester equation (cost $O(N^3)$), avoiding the $N^2\!\times\!N^2$
Jacobian entirely.  This method is L-stable and unconditionally stable
at $h\cdot\lambda_{\max}$ up to 181; at $K=200$ steps it achieves a
$14\times$ wall-time speedup over Magnus, with step count independent
of NQuad.  Richardson extrapolation raises the reflection operator to
second order at cost $3K$, further reducing the required $K$ by
${\sim}10\times$ for $10^{-3}$ accuracy; the overall convergence is
then limited by the first-order companion transmission equation, which
a proper two-stage Rosenbrock method would resolve.

These results establish the theoretical and numerical foundations for
diffusion-domain methods while honestly delineating their current
limitations for the problem sizes of interest.
\end{abstract}

\tableofcontents
\vspace{1em}

%======================================================================
\section{Introduction and Motivation}\label{sec:intro}
%======================================================================

\subsection{The stability ceiling}

Report~I~\cite{report1} established that the star-product Magnus
integrator requires $K\gtrsim\lambda_{\max}\tbot$ steps to maintain
per-step conditioning $\cond(\mat{a})\lesssim 1.3$.  Here
$\lambda_{\max}$ is the spectral radius of the DISORT coefficient
matrix $\mat{A}$, $h=\tbot/K$ is the step size, and the constraint is
$h\cdot\lambda_{\max}\lesssim 1$.  For NQuad${}=8$ ($N\!=\!4$),
$\lambda_{\max}\approx 14$; for NQuad${}=16$ ($N\!=\!8$),
$\lambda_{\max}\approx 45$.

For thick, highly scattering atmospheres---the diffusion domain---the
physics is governed by the \emph{diffusion eigenvalue} $\kmin$, which
is much smaller than $\lambda_{\max}$.  The step-count constraint
$K\propto\lambda_{\max}\tbot$ forces the solver to resolve
eigenmodes that decay on a scale $\sim 1/\lambda_{\max}$, even though
the observable reflectance and flux change on the much longer scale
$\sim 1/\kmin$.

\begin{remark}[Numerical example]
For $\omega=0.999$, $g=0.85$, NQuad${}=16$: $\lambda_{\max}\approx 45$,
$\kmin\approx 0.021$, gap $\gamma\approx 2128$.  A slab of
$\tsub=20$ requires $K\approx 913$ Magnus steps.  If the solver
could step at the diffusion rate, $K\sim\kmin\cdot\tsub\approx 0.4$
would suffice for accuracy---a factor of ${\sim}2000$ fewer steps.
\end{remark}

\subsection{Context: two failed approaches}

Before the star-product formulation of Report~I, two approaches to
thick-atmosphere integration were attempted and failed:
\begin{enumerate}[nosep]
\item \textbf{Step-by-step SVD propagator}: attempted to track
  singular values of the $2N\!\times\!2N$ propagator at each step to
  separate growing and decaying modes.  Failed because the DISORT
  matrix $\mat{A}$ is non-normal, causing the per-step SVD to diverge
  from the true singular values for $\tau\gtrsim 1.65$.
\item \textbf{Frozen-mode propagator}: attempted to freeze the mode
  structure from an initial eigendecomposition and propagate in the
  eigenbasis.  Failed for the same fundamental reason: non-normality
  of $\mat{A}$ means its eigenvectors are not orthogonal, and
  small perturbations in the fast-decaying modes contaminate the
  slow modes.
\end{enumerate}
Both approaches operated on the $2N\!\times\!2N$ propagator, where
exponentially growing and decaying modes coexist.  The star-product
formulation eliminates this by working entirely with $O(1)$-bounded
$N\!\times\!N$ operators.  The question is whether the
\emph{step-count} constraint can also be relaxed.

\subsection{Outline}

\Cref{sec:eigenvalue_gap} quantifies the eigenvalue gap across
parameter space.  \Cref{sec:riccati} derives the matrix Riccati
equation from invariant imbedding.  \Cref{sec:R_inf} develops the
semi-infinite reflectance and its computation via iterative doubling.
\Cref{sec:companion} presents the companion equations for
transmission and source.  \Cref{sec:decomposition} formulates the
three-domain star-product decomposition.  \Cref{sec:safety} addresses
why this approach is safe compared to the failed alternatives.
\Cref{sec:solvers} presents the numerical benchmarks comparing
explicit and implicit solvers.  \Cref{sec:open} discusses open
questions.


%======================================================================
\section{Eigenvalue Structure and the Diffusion Gap}
\label{sec:eigenvalue_gap}
%======================================================================

\subsection{Eigenvalue reduction}

The $2N\!\times\!2N$ coefficient matrix has the block structure
(Report~I, eq.~(5)):
\begin{equation}\label{eq:A_blocks}
  \mat{A} = \begin{pmatrix}-\mat{\alpha} & -\mat{\beta}\\
  \mat{\beta} & \mat{\alpha}\end{pmatrix},
\end{equation}
where $\mat{\alpha},\mat{\beta}\in\R^{N\times N}$ depend on $\omega$,
the phase function, and the quadrature.  The eigenvalues of $\mat{A}$
come in $\pm k_j$ pairs.  Defining sum and difference variables
$\vect{v}=\vect{u}^+ + \vect{u}^-$ and $\vect{w}=\vect{u}^+ - \vect{u}^-$,
the $2N$-dimensional eigenvalue problem decouples into two
$N$-dimensional problems:
\begin{equation}\label{eq:k_reduction}
  k_j^2 = \text{eigenvalues of}\;\;
  (\mat{\alpha}-\mat{\beta})(\mat{\alpha}+\mat{\beta}).
\end{equation}
All $k_j^2$ are real and non-negative for physical $\omega\in[0,1)$.
We order $k_1 \geq k_2 \geq \cdots \geq k_N > 0$ and define:
\begin{itemize}[nosep]
\item $\kmax = k_1$: the \emph{ballistic} eigenvalue (fastest mode),
\item $\kmin = k_N$: the \emph{diffusion} eigenvalue (slowest mode),
\item $\gamma = \kmax/\kmin$: the \emph{eigenvalue gap ratio}.
\end{itemize}

\subsection{Asymptotic behaviour}

The ballistic eigenvalue scales with the quadrature:
\begin{equation}
  \kmax \approx \frac{1}{\mu_{\min}} \quad\text{(weakly dependent on $\omega$, $g$)},
\end{equation}
where $\mu_{\min}$ is the smallest quadrature abscissa.  Since
$\mu_{\min}\sim O(N^{-2})$ for Gauss--Legendre quadrature, $\kmax$
grows as $O(N^2)$.

The diffusion eigenvalue, by contrast, is a \emph{physical} quantity
that converges rapidly with $N$.  The two-stream approximation gives
the closed-form prediction:
\begin{equation}\label{eq:k_2s}
  k_{2\mathrm{s}} = \sqrt{3(1-\omega)(1-\omega g)},
\end{equation}
and the multi-stream $\kmin$ satisfies $\kmin/k_{2\mathrm{s}}\to 1$ as
$\omega\to 1$ (the diffusion limit).  \Cref{tab:eigenvalues} confirms
this numerically.

\begin{table}[h]\centering
\caption{Eigenvalue gap for representative atmospheres (NQuad${}=8$,
$N\!=\!4$, $g=0.85$).  $k_{2\mathrm{s}}$ is the two-stream prediction~\eqref{eq:k_2s}.}
\label{tab:eigenvalues}
\begin{tabular}{cccccc}
\toprule
$\omega$ & $\kmax$ & $\kmin$ & $\gamma$ & $k_{2\mathrm{s}}$ & $\kmin/k_{2\mathrm{s}}$ \\
\midrule
0.50   & 12.40 & 0.661 &  18.8 & 0.929 & 0.712 \\
0.90   & 10.68 & 0.237 &  45.1 & 0.266 & 0.893 \\
0.99   & 10.28 & 0.068 & 151.1 & 0.069 & 0.986 \\
0.999  & 10.23 & 0.021 & 481.8 & 0.021 & 0.999 \\
0.9999 & 10.23 & 0.007 & 1524.9 & 0.007 & 1.000 \\
\bottomrule
\end{tabular}
\end{table}

\subsection{NQuad dependence}

The gap ratio grows as $\gamma\sim O(N^2)$ because $\kmax$ grows with
$1/\mu_{\min}$ while $\kmin$ is approximately constant.
\Cref{tab:nquad_gap} shows this for $\omega=0.99$, $g=0.85$.

\begin{table}[h]\centering
\caption{NQuad dependence of the eigenvalue gap ($\omega=0.99$,
$g=0.85$).}
\label{tab:nquad_gap}
\begin{tabular}{ccccc}
\toprule
NQuad & $N$ & $\kmax$ & $\kmin$ & $\gamma$ \\
\midrule
 4 &  2 &   1.82 & 0.068 &   26.9 \\
 8 &  4 &  10.28 & 0.068 &  151.1 \\
16 &  8 &  45.24 & 0.068 &  665.4 \\
\bottomrule
\end{tabular}
\end{table}

For operational retrievals with NQuad${}=32$ ($N\!=\!16$):
$\kmax\sim 250$, $\gamma\sim 3700$ for $\omega=0.99$, and
$\gamma\sim 37000$ for $\omega=0.999$.  The step-count
$K\gtrsim\lambda_{\max}\tbot$ becomes prohibitive for thick clouds.

\subsection{Variation along a cloud profile}

For an adiabatic cloud profile (linear $\omega(\tau)$ and $g(\tau)$),
the eigenvalue gap varies smoothly: from $\gamma\approx 37$ at cloud top
($\omega=0.85$) to $\gamma\approx 69$ at cloud base ($\omega=0.96$)
for NQuad${}=8$.  The smoothness of $\gamma(\tau)$ supports the idea
that a diffusion-domain solver can track the slowly varying physics
without resolving the fast modes at every step.


%======================================================================
\section{The Matrix Riccati Equation}\label{sec:riccati}
%======================================================================

\subsection{Derivation from invariant imbedding}

The matrix Riccati equation describes how the reflection operator
$\mat{R}(\sigma)$ of a slab of thickness $\sigma$ evolves as
the slab is built up from below, one infinitesimal layer at a time.
This is the \emph{invariant imbedding} formulation of
Bellman, Kalaba, and Prestrud~\cite{bellman1960}.

Consider a slab of thickness $\sigma$ with reflection operator
$\mat{R}(\sigma)$, and append an infinitesimal layer of thickness
$d\sigma$ at the bottom.  The infinitesimal layer has:
\begin{align}
  \mat{r}_\uparrow &= \mat{\beta}\dd\sigma + O(d\sigma^2), &
  \mat{t}_\uparrow &= \mat{I} + \mat{\alpha}\dd\sigma + O(d\sigma^2), \\
  \mat{r}_\downarrow &= \mat{\beta}\dd\sigma + O(d\sigma^2), &
  \mat{t}_\downarrow &= \mat{I} + \mat{\alpha}\dd\sigma + O(d\sigma^2),
\end{align}
where $\mat{\alpha}$ and $\mat{\beta}$ are the absorption--scattering
blocks of \eqref{eq:A_blocks}.

Applying the star product (Report~I, Definition~2) to combine the
existing slab (on top) with the infinitesimal layer (below) and
retaining terms to first order in $d\sigma$:
\begin{equation}\label{eq:riccati_derivation}
  \mat{R}(\sigma+d\sigma)
  = \mat{R} + \bigl(\mat{\alpha}\mat{R} + \mat{R}\mat{\alpha}
    + \mat{R}\mat{\beta}\mat{R} + \mat{\beta}\bigr)\dd\sigma
    + O(d\sigma^2).
\end{equation}
This yields the \textbf{matrix Riccati ODE}:
\begin{equation}\label{eq:riccati_pos}
  \boxed{
  \frac{d\mat{R}}{d\sigma}
  = \mat{\alpha}\mat{R} + \mat{R}\mat{\alpha}
    + \mat{R}\mat{\beta}\mat{R} + \mat{\beta},
  \qquad \mat{R}(0)=\mat{0}.
  }
\end{equation}
Here $\sigma$ is the slab thickness increasing from $0$ to $\tsub$.

\begin{remark}[Sign convention]\label{rem:sign_convention}
The sign in~\eqref{eq:riccati_pos} is \emph{positive} because the slab
is built from below: $\sigma$ increases as material is added at the
bottom.  Report~I, equation~(16), presents the same equation with
\emph{negative} signs; that convention describes $\mat{R}(\tau)$ as a
function of the optical-depth coordinate $\tau$ increasing
\emph{downward into the slab from the top}.  The two forms are related
by $\sigma=\tsub-\tau$ (i.e., $d\sigma = -d\tau$), and both describe
the same physical process.  For \emph{numerical integration}, the
positive-sign form~\eqref{eq:riccati_pos} is the natural choice: one
integrates forward from $\sigma=0$ (empty slab) to $\sigma=\tsub$
(full sub-domain), with $\mat{R}(\sigma)$ monotonically increasing in
norm from $\mat{0}$ toward $\Rinf$.
\end{remark}

\subsection{Connection to the M\"obius transformation}

The exact one-step Riccati update is the M\"obius transformation
(Report~I, equation~(17)):
\begin{equation}
  \mat{R}(\sigma+h)
  = (\mat{a}\,\mat{R}(\sigma)+\mat{b})\,
    (\mat{c}\,\mat{R}(\sigma)+\mat{d})^{-1},
\end{equation}
where $\bigl[\begin{smallmatrix}\mat{a}&\mat{b}\\
\mat{c}&\mat{d}\end{smallmatrix}\bigr]=\exp(h\mat{A})$.
The Magnus integrator computes this via \texttt{expm} and the
star product, which is equivalent to applying the M\"obius
transformation at each step.  The Riccati ODE~\eqref{eq:riccati_pos}
is the infinitesimal generator of this M\"obius semigroup.

\subsection{Linearisation and stiffness}\label{sec:linearisation}

As $\sigma\to\infty$, the reflection operator converges to the
semi-infinite reflectance $\Rinf$ (\Cref{sec:R_inf}).  Linearising
around $\Rinf$ with $\mat{R}=\Rinf+\mat{E}$:
\begin{equation}\label{eq:linearised}
  \frac{d\mat{E}}{d\sigma}
  = (\mat{\alpha}+\Rinf\mat{\beta})\mat{E}
  + \mat{E}(\mat{\alpha}+\mat{\beta}\Rinf)
  + \mat{E}\mat{\beta}\mat{E}.
\end{equation}
Dropping the quadratic term and vectorising $\vect{e}=\operatorname{vec}(\mat{E})$
(row-major, following numpy's C-order \texttt{ravel}):
\begin{equation}\label{eq:linearised_vec}
  \frac{d\vect{e}}{d\sigma}
  = \bigl[
    (\mat{\alpha}+\Rinf\mat{\beta})\otimes\mat{I}_N
    + \mat{I}_N\otimes(\mat{\alpha}+\mat{\beta}\Rinf)^T
  \bigr]\,\vect{e},
\end{equation}
where $\otimes$ denotes the Kronecker product.  The eigenvalues of
the coefficient matrix in~\eqref{eq:linearised_vec} are
$\{k_i+k_j : 1\leq i,j\leq N\}$, all positive (since
each $k_j>0$).

\begin{proposition}[Stiffness of the Riccati ODE]\label{prop:stiffness}
The Riccati ODE~\eqref{eq:riccati_pos} has the following
eigenvalue structure near $\Rinf$:
\begin{enumerate}[nosep]
\item Slowest mode: eigenvalue $2\kmin$ (diffusion--diffusion
  interaction);
\item Fastest mode: eigenvalue $2\kmax$ (ballistic--ballistic
  interaction);
\item Stiffness ratio: $\kmax/\kmin = \gamma$, the same eigenvalue gap
  as the original coefficient matrix.
\end{enumerate}
The Riccati ODE is therefore \emph{stiff} with stiffness ratio $\gamma$.
For $\omega=0.99$, $g=0.85$, NQuad${}=16$: $\gamma\approx 665$.
\end{proposition}

\begin{remark}[Critical consequence for explicit integrators]
An explicit integrator (e.g., RK4) applied to~\eqref{eq:riccati_pos}
must satisfy the CFL-type constraint $h\cdot 2\kmax \lesssim C$ for
stability, where $C\approx 2.8$ for RK4.  This gives
$K\gtrsim 2\kmax\tsub/C$---the \emph{same} step constraint as the
Magnus integrator.  The Riccati reformulation does not improve the
situation for explicit methods.
\end{remark}

\subsection{Jacobian for implicit solvers}

Vectorising the full nonlinear Riccati
ODE~\eqref{eq:riccati_pos} with
$\vect{r}=\operatorname{vec}(\mat{R})$ (row-major):
\begin{equation}\label{eq:riccati_vec}
  \frac{d\vect{r}}{d\sigma}
  = \operatorname{vec}\bigl(
    \mat{\alpha}\mat{R} + \mat{R}\mat{\alpha}
    + \mat{R}\mat{\beta}\mat{R} + \mat{\beta}
  \bigr).
\end{equation}
The Jacobian (required by implicit solvers such as Radau~IIA) is:
\begin{equation}\label{eq:jac_riccati}
  \mat{J}_R
  = (\mat{\alpha}+\mat{R}\mat{\beta})\otimes\mat{I}_N
  + \mat{I}_N\otimes(\mat{\alpha}+\mat{\beta}\mat{R})^T.
\end{equation}
This is an $N^2\!\times\!N^2$ matrix.  For $N=4$ (NQuad${}=8$):
$16\times 16$, trivially cheap.  For $N=8$ (NQuad${}=16$):
$64\times 64$, still manageable.  For $N=16$ (NQuad${}=32$):
$256\times 256$, feasible but not negligible.


%======================================================================
\section{Semi-Infinite Reflectance}\label{sec:R_inf}
%======================================================================

\subsection{The algebraic Riccati equation}

As $\sigma\to\infty$, $d\mat{R}/d\sigma\to\mat{0}$, and
$\mat{R}\to\Rinf$ satisfying the \emph{algebraic Riccati equation}:
\begin{equation}\label{eq:ARE}
  \mat{\alpha}\Rinf + \Rinf\mat{\alpha}
  + \Rinf\mat{\beta}\Rinf + \mat{\beta} = \mat{0}.
\end{equation}
For a homogeneous atmosphere, $\Rinf$ is the reflectance of a
semi-infinite slab.  It exists and is unique for $\omega\in(0,1)$
and $g\in[0,1)$~\cite{bittanti1991,lancaster1995}.  Its spectral norm satisfies
$\|\Rinf\|_2 < 1$ (energy conservation: some energy is absorbed at
each scattering event in the infinite slab).

\subsection{Computation via iterative doubling}

$\Rinf$ can be computed without solving the
ARE~\eqref{eq:ARE} directly, by exploiting the star product.
Start with the operators $(\mat{R},\mat{T}_\uparrow,\mat{T}_\downarrow,\mat{R}_\downarrow,\vect{s}_\uparrow,\vect{s}_\downarrow)$ for a finite slab of thickness $\tau_0$
(computed by the Magnus star-product algorithm of Report~I).  Then
\emph{self-double}: combine the slab with itself via the star product.
After $s$ doublings, the slab has thickness $\tau_0\cdot 2^s$, and
$\mat{R}_\uparrow\to\Rinf$ exponentially fast.

\Cref{tab:doubling} shows convergence for three representative cases.

\begin{table}[h]\centering
\caption{Iterative doubling convergence to $\Rinf$ (NQuad${}=8$,
initial $\tau_0=1$ with 200 Magnus steps).  $\|\Delta\mat{R}\|$ is
$\|\mat{R}_{2\tau}-\mat{R}_\tau\|_2$.}\label{tab:doubling}
\begin{tabular}{lccc}
\toprule
$\tau$ & $\omega=0.9,\;g=0.5$ & $\omega=0.99,\;g=0.85$ & $\omega=0.999,\;g=0.85$ \\
\midrule
  2  & $8.1\times 10^{-2}$ & $9.5\times 10^{-2}$ & $1.0\times 10^{-1}$ \\
  4  & $5.0\times 10^{-2}$ & $1.1\times 10^{-1}$ & $1.3\times 10^{-1}$ \\
  8  & $1.3\times 10^{-2}$ & $1.3\times 10^{-1}$ & $1.6\times 10^{-1}$ \\
 16  & $6.0\times 10^{-4}$ & $9.9\times 10^{-2}$ & $1.6\times 10^{-1}$ \\
 32  & $1.3\times 10^{-6}$ & $4.1\times 10^{-2}$ & $1.3\times 10^{-1}$ \\
 64  & $5.5\times 10^{-12}$ & $5.0\times 10^{-3}$ & $6.7\times 10^{-2}$ \\
128  & $< 10^{-15}$ & $6.4\times 10^{-5}$ & $1.8\times 10^{-2}$ \\
256  & ---          & $1.1\times 10^{-8}$ & $1.2\times 10^{-3}$ \\
\bottomrule
\end{tabular}
\end{table}

As expected, the convergence slows as $\omega\to 1$ ($\kmin\to 0$),
since the number of doublings required scales as
$s\gtrsim \log_2(1/\kmin)$.

\subsection{Convergence rate}

\begin{proposition}[Convergence of $\mat{R}(\tau)$ to $\Rinf$]
\label{prop:convergence_R_inf}
For a homogeneous slab:
\begin{equation}
  \|\mat{R}(\tau) - \Rinf\|_2 \leq C\,e^{-2\kmin\tau},
\end{equation}
where $C$ depends on the initial mismatch $\|\mat{R}(0)-\Rinf\|$.
The rate $2\kmin$ follows from the slowest eigenvalue of the
linearisation~\eqref{eq:linearised_vec}.
\end{proposition}

Numerically, for $\omega=0.99$, $g=0.85$ ($\kmin=0.068$):
\begin{center}
\begin{tabular}{ccc}
\toprule
$\tau$ & $\|\mat{R}(\tau)-\Rinf\|_2$ & $e^{-2\kmin\tau}$ \\
\midrule
0.5  & $5.3\times 10^{-1}$ & $9.3\times 10^{-1}$ \\
5.0  & $2.3\times 10^{-1}$ & $5.1\times 10^{-1}$ \\
20.0 & $2.6\times 10^{-2}$ & $6.6\times 10^{-2}$ \\
50.0 & $4.3\times 10^{-4}$ & $1.1\times 10^{-3}$ \\
\bottomrule
\end{tabular}
\end{center}
The ratio $\|\mat{R}-\Rinf\|/e^{-2\kmin\tau}$ is approximately
constant (${\approx}0.4$), confirming the predicted decay rate.

\subsection{$\Rinf$ for $\tau$-varying atmospheres}

For a $\tau$-varying atmosphere, $\mat{\alpha}(\tau)$ and
$\mat{\beta}(\tau)$ change with depth, and the global $\Rinf$ does not
exist.  However, one can define the \emph{local} $\Rinf(\tau)$: the
solution to~\eqref{eq:ARE} using the frozen local values of
$\mat{\alpha}(\tau)$, $\mat{\beta}(\tau)$.

Along an adiabatic cloud profile ($\omega$ from 0.85 to 0.96, $g$ from
0.865 to 0.820 over $\tbot=30$), the local $\|\Rinf(\tau)\|_2$
varies smoothly from 0.27 to 0.47, with
$\|\Rinf(\tau_j)-\Rinf(\tau_{j-1})\|\lesssim 0.04$.  This
supports the physical picture of \emph{adiabatic tracking}: in the
diffusion domain, $\mat{R}(\sigma)$ approximately tracks the slowly
varying $\Rinf(\tau)$ without needing to resolve the fast transient
modes.


%======================================================================
\section{Companion Equations for T and s}
\label{sec:companion}
%======================================================================

The reflection operator $\mat{R}$ is only one of the six slab
operators needed for the star-product coupling.  The transmission
operator $\mat{T}_\uparrow$ and source vector $\vect{s}_\uparrow$
satisfy companion ODEs that are \emph{linear} given $\mat{R}(\sigma)$.

\subsection{Transmission equation}

\begin{proposition}[Companion T equation]\label{prop:companion_T}
The upward transmission operator satisfies
\begin{equation}\label{eq:T_ode}
  \frac{d\mat{T}}{d\sigma} = \mat{T}\,(\mat{\alpha}+\mat{\beta}\mat{R}),
  \qquad \mat{T}(0) = \mat{I}_N.
\end{equation}
\end{proposition}

\begin{proof}[Derivation]
The star product for combining the accumulated slab with an
infinitesimal layer gives, to first order:
\begin{equation}
  \mat{T}_\uparrow(\sigma+d\sigma)
  = \mat{T}_\uparrow\,(\mat{I}+\mat{\alpha}\dd\sigma)
    \,(\mat{I}-\mat{\beta}\dd\sigma\cdot\mat{R}_\downarrow)^{-1}
    \cdot\mat{t}_\uparrow.
\end{equation}
Expanding and using $\mat{R}_\downarrow\approx\mat{R}_\uparrow$
(which holds for homogeneous slabs by symmetry of the $\mat{A}$
block structure) and $\mat{t}_\uparrow=\mat{I}+\mat{\alpha}\dd\sigma$:
\begin{equation}
  d\mat{T} = \mat{T}\,(\mat{\alpha}+\mat{\beta}\mat{R})\dd\sigma.
\end{equation}
\end{proof}

\begin{remark}[Right multiplication]
The right multiplication in~\eqref{eq:T_ode} (as opposed to left
multiplication $(\mat{\alpha}+\mat{\beta}\mat{R})\mat{T}$) is
essential.  It follows from the star-product structure: the
accumulated transmission $\mat{T}_\uparrow$ acts first, then the
infinitesimal layer modifies the result.
\end{remark}

Vectorising $\vect{t}=\operatorname{vec}(\mat{T})$ (row-major):
\begin{equation}\label{eq:T_vec}
  \frac{d\vect{t}}{d\sigma}
  = \bigl[\mat{I}_N\otimes(\mat{\alpha}+\mat{\beta}\mat{R})^T\bigr]\,\vect{t}.
\end{equation}
The Jacobian is the $N^2\!\times\!N^2$ matrix
$\mat{J}_T = \mat{I}_N\otimes(\mat{\alpha}+\mat{\beta}\mat{R})^T$,
which has the same eigenvalues $\{k_j : j=1,\ldots,N\}$ (each with
multiplicity $N$).  The stiffness ratio is again $\kmax/\kmin=\gamma$.

\subsection{Source equation}

For completeness, the source vector satisfies the inhomogeneous ODE:
\begin{equation}\label{eq:s_ode}
  \frac{d\vect{s}_\uparrow}{d\sigma}
  = (\mat{\alpha}+\mat{R}\mat{\beta})\,\vect{s}_\uparrow
  + \text{(beam source terms)},
  \qquad \vect{s}_\uparrow(0) = \vect{0}.
\end{equation}
This is an $N$-dimensional linear ODE (given $\mat{R}(\sigma)$) with
the same stiffness structure.

\subsection{Total degrees of freedom}

The full Riccati + companion system has:
\begin{itemize}[nosep]
\item $N^2$ DOFs for $\mat{R}$ (symmetric: could be reduced to
  $N(N+1)/2$, but this complicates the Jacobian),
\item $N^2$ DOFs for $\mat{T}_\uparrow$,
\item $N$ DOFs for $\vect{s}_\uparrow$,
\end{itemize}
totalling $2N^2+N$.  Compare this with $4N^2$ for the
$2N\!\times\!2N$ propagator.  For $N=8$: $136$ vs $256$ DOFs.


%======================================================================
\section{Domain Decomposition via Star Product}
\label{sec:decomposition}
%======================================================================

\subsection{Three-domain structure}

We decompose the atmosphere $[0,\tbot]$ into three domains:
\begin{center}
\begin{tabular}{lll}
\textbf{Domain~I}   & $[0, \tau_1]$         & Non-diffusive (beam active) \\
\textbf{Domain~II}  & $[\tau_1, \tau_2]$    & Diffusion domain \\
\textbf{Domain~III} & $[\tau_2, \tbot]$     & Non-diffusive (near surface)
\end{tabular}
\end{center}
Domain~II is the optically thick, highly scattering interior where the
direct beam is negligible ($e^{-\tau_1/\mu_0}\ll 1$) and the physics
is dominated by the diffusion eigenvalue $\kmin$.  The boundaries
$\tau_1$ and $\tau_2$ are chosen so that the beam attenuation at
$\tau_1$ is negligible (e.g., $e^{-\tau_1/\mu_0}<10^{-8}$).

\subsection{Independent solution of each domain}

Each domain is solved independently, producing its own 6-tuple of
$N\!\times\!N$ scattering operators:
\begin{align}
  \text{Domain~I:}\quad  &(\mat{R}_{\uparrow}^{\mathrm{I}},\;\mat{T}_{\uparrow}^{\mathrm{I}},\;\mat{T}_{\downarrow}^{\mathrm{I}},\;\mat{R}_{\downarrow}^{\mathrm{I}},\;\vect{s}_{\uparrow}^{\mathrm{I}},\;\vect{s}_{\downarrow}^{\mathrm{I}}), \\
  \text{Domain~II:}\quad &(\mat{R}_{\uparrow}^{\mathrm{II}},\;\mat{T}_{\uparrow}^{\mathrm{II}},\;\mat{T}_{\downarrow}^{\mathrm{II}},\;\mat{R}_{\downarrow}^{\mathrm{II}},\;\vect{s}_{\uparrow}^{\mathrm{II}},\;\vect{s}_{\downarrow}^{\mathrm{II}}), \\
  \text{Domain~III:}\quad &(\mat{R}_{\uparrow}^{\mathrm{III}},\;\mat{T}_{\uparrow}^{\mathrm{III}},\;\mat{T}_{\downarrow}^{\mathrm{III}},\;\mat{R}_{\downarrow}^{\mathrm{III}},\;\vect{s}_{\uparrow}^{\mathrm{III}},\;\vect{s}_{\downarrow}^{\mathrm{III}}).
\end{align}
The choice of solver for each domain is arbitrary:
\begin{itemize}[nosep]
\item Domain~I: full Magnus star product (beam source is active),
\item Domain~II: implicit Riccati solver \emph{or} Magnus star product,
\item Domain~III: full Magnus star product.
\end{itemize}

\subsection{Coupling via star product}

The three domains are coupled by two star-product operations:
\begin{equation}
  (\text{I}\star\text{II})\star\text{III}
  = \text{I}\star(\text{II}\star\text{III})
  \;\longrightarrow\;
  (\mat{R}_\uparrow,\mat{T}_\uparrow,\mat{T}_\downarrow,\mat{R}_\downarrow,\vect{s}_\uparrow,\vect{s}_\downarrow)_{\text{total}}.
\end{equation}
The star product is associative (Report~I, Definition~2),
so the order does not matter.  The total operators then enter the
standard BC solver (Report~I, \S6).

\subsection{Numerical validation}

\Cref{tab:3domain} compares the three-domain decomposition with a
single full-atmosphere Magnus solve for $\tbot=50$, $\omega=0.99$,
$g=0.85$, NQuad${}=16$, $\mu_0=0.5$, $I_0=1$, with $\tau_1=10$,
$\tau_2=30$.

\begin{table}[h]\centering
\caption{Three-domain decomposition vs full Magnus
($\tbot=50$, $\omega=0.99$, $g=0.85$, NQuad${}=16$).}
\label{tab:3domain}
\begin{tabular}{lcc}
\toprule
& Full Magnus & 3-Domain \\
\midrule
Total steps $K$ & 2272 & 462+914+914=2290 \\
Flux at ToA & $2.9957\times 10^{-1}$ & $2.9957\times 10^{-1}$ \\
Rel.\ err.\ vs pydisort & $9.0\times 10^{-5}$ & $8.7\times 10^{-5}$ \\
$\|\vect{u}_{\text{3d}}-\vect{u}_{\text{full}}\|/\|\vect{u}_{\text{full}}\|$ & --- & $2.1\times 10^{-5}$ \\
Wall time (s) & 2.69 & 2.71 \\
\bottomrule
\end{tabular}
\end{table}

The three-domain coupling introduces no additional error beyond the
$O(h^2)$ discretisation error of each domain's solver.  The star
product is exact---it couples the domains without any approximation.


%======================================================================
\section{Why This Approach Is Safe}\label{sec:safety}
%======================================================================

The two previous approaches to thick atmospheres (SVD propagator,
frozen modes) failed catastrophically.  It is important to explain
why the Riccati/star-product domain decomposition is fundamentally
different, and where it is \emph{not} different.

\subsection{What failed before}

Both previous approaches operated on the $2N\!\times\!2N$ propagator
$\mat{\Phi}(\tau)$, which contains both growing and decaying modes.
For thick slabs, $\cond(\mat{\Phi})=e^{2\lambda_{\max}\tau}\gg 1$.
Any numerical error in the growing modes contaminates the decaying
modes (and vice versa), because $\mat{A}$ is non-normal and its
eigenvectors are not orthogonal.

\begin{itemize}[nosep]
\item \textbf{SVD approach}: tracked singular values per step, hoping
  to separate fast/slow modes.  But the per-step SVD singular values
  (which are all near~1 for small $h$) bear no relation to the
  cumulative singular values (which span $e^{\pm\lambda_{\max}\tau}$).
  The approach diverged for $\tau\gtrsim 1.65$.
\item \textbf{Frozen-mode approach}: diagonalised $\mat{A}$ once and
  propagated in the eigenbasis.  But the non-orthogonality of the
  eigenvectors means that projecting onto ``slow modes'' while
  discarding ``fast modes'' is ill-conditioned.  Small errors in the
  projection grew exponentially.
\end{itemize}

\subsection{Why $N\!\times\!N$ operators are safe}

The Riccati/star-product formulation works entirely with $N\!\times\!N$
operators: $\mat{R}$, $\mat{T}$, $\vect{s}$.  All of these are
bounded by energy conservation (Proposition~1 of Report~I):
\begin{itemize}[nosep]
\item $\|\mat{R}\|_2 < 1$ (reflection is sub-unitary for $\omega<1$),
\item $\|\mat{T}\|_2 \leq 1$ (transmission cannot exceed unity),
\item $\|\vect{s}\|$ is $O(1)$ (bounded by the source intensity).
\end{itemize}
No exponentially growing quantities appear anywhere.  The star product
preserves these bounds at each step.

\subsection{Error containment}

\begin{proposition}[Error containment in the star product]
If a diffusion-domain solver produces $\mat{R}^{\mathrm{II}}$ with
error $\|\mat{R}^{\mathrm{II}}-\mat{R}^{\mathrm{II}}_{\text{exact}}\|
= \epsilon$, then the error in the total coupled operators is
$O(\epsilon)$---not $O(e^{\kmax\tau}\epsilon)$.
\end{proposition}

This follows from the star product being a \emph{contraction}:
bounded inputs produce bounded outputs, with Lipschitz constant
depending on $\|\mat{R}\|_2<1$ and $\cond(\mat{I}-\mat{r}\mat{R})$.
The worst case is a modest error amplification factor, not an
exponential blowup.

\subsection{What is NOT different}

The Riccati ODE~\eqref{eq:riccati_pos} is stiff with the same
eigenvalue gap $\gamma$ as the original problem.  An explicit
integrator applied to the Riccati ODE has the same step-count
constraint as the Magnus integrator.  The improvement comes only
from using an \emph{implicit} solver that can handle stiff
systems---and even then, the improvement is in the number of
\emph{function evaluations}, not necessarily in wall time.

The critical safety guarantee is: if the diffusion-domain solver
fails or is inaccurate, the errors remain $O(\epsilon)$, not
$O(e^{\kmax\tau}\epsilon)$.  One can always fall back to the full
Magnus star product---the worst case is the same step count, not
a catastrophic failure.


%======================================================================
\section{Candidate Solvers and Stiffness Analysis}\label{sec:solvers}
%======================================================================

This section presents the numerical benchmark comparing three solvers
on the diffusion sub-domain $[\tau_1,\tau_2]=[10,30]$ of a thick slab
($\tbot=50$, $\omega=0.99$, $g=0.85$, NQuad${}=16$, $N=8$).  The
sub-domain boundary conditions $\vect{b}_{\text{neg}}$ (downwelling
at $\tau_1$) and $\vect{b}_{\text{pos}}$ (upwelling at $\tau_2$) are
extracted from the exact pydisort solution, providing a known reference.

\subsection{Explicit RK4 on the Riccati ODE}

\Cref{tab:rk4} shows the behaviour of classical RK4 applied
to~\eqref{eq:riccati_pos}.

\begin{table}[h]\centering
\caption{RK4 stability on the Riccati ODE ($\tsub=20$,
$\kmax=45.2$, NQuad${}=16$).}\label{tab:rk4}
\begin{tabular}{cccc}
\toprule
$K$ & $h\cdot\kmax$ & Status & Rel.\ err. \\
\midrule
50  & 18.1  & blowup & --- \\
100 & 9.0   & blowup & --- \\
200 & 4.5   & blowup & --- \\
400 & 2.3   & blowup & --- \\
600 & 1.5   & OK & $3.4\times 10^{-3}$ \\
800 & 1.1   & OK & $4.6\times 10^{-5}$ \\
1500 & 0.6  & OK & $9.8\times 10^{-6}$ \\
\bottomrule
\end{tabular}
\end{table}

The stability boundary is at $h\cdot\kmax\approx 1.5$, consistent
with the RK4 stability region applied to the fastest eigenvalue
$2\kmax$.  The minimum step count ($K\approx 600$) is comparable
to the Magnus step count ($K=914$).

\textbf{Critical negative result}: explicit integration of the
Riccati ODE offers no speedup over the Magnus star product.  Both
are constrained by the same ballistic eigenvalue $\kmax$.

\subsection{Implicit Radau~IIA on the Riccati ODE}

The implicit Radau~IIA solver (\texttt{scipy.integrate.solve\_ivp}
with \texttt{method='Radau'}) handles stiff systems by allowing the
step size to be determined by accuracy rather than stability.  The
$N^2=64$ DOFs for $\mat{R}$ are vectorised, and the analytical
Jacobian~\eqref{eq:jac_riccati} is provided.

After solving $\mat{R}(\sigma)$, the companion $\mat{T}(\sigma)$
equation~\eqref{eq:T_vec} is integrated as a separate linear ODE (also
via Radau, with its own Jacobian).

Results:
\begin{itemize}[nosep]
\item Riccati $\mat{R}$: 1925 function evaluations, 4 Jacobian
  evaluations, 84 LU factorisations.
\item Companion $\mat{T}$: 2895 function evaluations.
\item Total: $\approx 4820$ function evaluations.
\item Accuracy: $\|\vect{u}_{\text{Radau}}-\vect{u}_{\text{ref}}\|/\|\vect{u}_{\text{ref}}\| = 5.1\times 10^{-9}$.
\item Wall time: 0.34~s (vs 0.68~s for Magnus).
\end{itemize}

The accuracy matches the Magnus solver to within discretisation error.
The wall-time advantage is modest (${\sim}2\times$).

\subsection{Scaling with $\omega$}\label{sec:omega_scaling}

The key question is whether the implicit solver's advantage grows
with the eigenvalue gap $\gamma$ (i.e., as $\omega\to 1$).
\Cref{tab:scaling} presents the results.

\begin{table}[h]\centering
\caption{Scaling with $\omega$ for the diffusion sub-domain
($\tsub=20$, $g=0.85$, NQuad${}=16$).  $K_{\text{M}}$ is the Magnus
step count, $n_{\text{R}}$ is the total Radau function evaluations.}
\label{tab:scaling}
\begin{tabular}{ccccccc}
\toprule
$\omega$ & $\kmin$ & $\gamma$ & $K_{\text{M}}$ & $n_{\text{R}}$
& $t_{\text{M}}$ (s) & $t_{\text{R}}$ (s) \\
\midrule
0.9    & 0.237 &   193 & 924  & 5123 & 0.66 & 0.45 \\
0.99   & 0.068 &   665 & 914  & 4820 & 0.71 & 0.42 \\
0.999  & 0.021 &  2128 & 913  & 4770 & 0.65 & 0.39 \\
0.9999 & 0.007 &  6736 & 913  & 4770 & 0.71 & 0.37 \\
\bottomrule
\end{tabular}
\end{table}

\textbf{Key observation}: the Radau function evaluation count is
approximately \emph{constant} across four orders of magnitude of the
eigenvalue gap ($\gamma$ from 193 to 6736).  However, it does not
\emph{decrease} dramatically either.  The Radau solver requires
$n_{\text{R}}\approx 5000$ evaluations regardless of $\omega$,
while the Magnus solver requires $K_{\text{M}}\approx 914$ steps
(also roughly constant, since $\kmax$ depends weakly on $\omega$).

The Radau wall time is consistently ${\sim}1.5$--$2\times$ faster
than Magnus.  But the advantage does \emph{not} scale with $\gamma$
as originally hoped.

\subsection{Analysis}

There are two reasons the implicit solver does not achieve the
theoretically expected speedup of $O(\gamma)$:

\begin{enumerate}
\item \textbf{Jacobian cost}: the Radau solver must form and
  factorise the $64\times 64$ Jacobian at each step.  While the
  step count is small ($\sim 80$ LU factorisations), each
  factorisation involves $O(N^6)$ work (forming the $N^2\times N^2$
  Kronecker product and computing the LU decomposition).  For $N=8$
  this is manageable but not negligible.

\item \textbf{Companion T equation}: the $\mat{T}$ equation is
  integrated as a separate ODE after $\mat{R}$ is known.  This
  second integration adds $\sim 2900$ function evaluations.  A
  coupled solver integrating $\mat{R}$ and $\mat{T}$ simultaneously
  could potentially reduce this overhead.

\item \textbf{Constant $\kmax$}: the dominant cost in both the Magnus
  and Radau approaches is set by $\kmax$, which depends on the
  quadrature (NQuad) rather than the single-scattering albedo.  The
  implicit solver amortises the fast modes via implicit treatment, but
  the startup cost (initial Jacobian, establishing the step size) is
  comparable to the steady-state cost of the Magnus approach.
\end{enumerate}

\begin{remark}[Where the Riccati approach would shine]
The implicit Riccati solver would show a stronger advantage for
larger NQuad ($N\gg 8$), where $\kmax\sim O(N^2)$ pushes the Magnus
step count much higher while the Radau step count remains
approximately constant.  For NQuad${}=32$ ($N=16$): $\kmax\sim 250$,
and the Magnus solver would need $K\sim 5000$ steps for $\tsub=20$,
while the Radau solver should still need $\sim 5000$ function
evaluations.  However, the Jacobian is now $256\times 256$,
and its cost must be carefully assessed.
\end{remark}


\subsection{Exploiting the Sylvester structure}
\label{sec:sylvester}

The Radau solver's $N^2\!\times\!N^2$ Jacobian bottleneck can be
circumvented entirely by observing that the Riccati
Jacobian~\eqref{eq:jac_riccati} is a \emph{Kronecker sum}:
\begin{equation}
  \mat{J}_R
  = \underbrace{(\mat{\alpha}+\mat{R}\mat{\beta})}_{\mat{A}_n}
    \otimes\mat{I}_N
  + \mat{I}_N\otimes
    \underbrace{(\mat{\alpha}+\mat{\beta}\mat{R})^T}_{\mat{B}_n^T}.
\end{equation}
This means that each implicit step of a linearly-implicit (Rosenbrock)
method reduces to a \emph{Sylvester equation}---an
$N\!\times\!N$ problem, not an $N^2\!\times\!N^2$ one.

\subsubsection{Rosenbrock--Euler with Sylvester solve}

The linearly-implicit Euler method (Rosenbrock--Euler, $\gamma=1$)
applied to~\eqref{eq:riccati_vec} solves at each step:
\begin{equation}\label{eq:sylvester_step}
  (\mat{I}_N - h\mat{A}_n)\,\Delta\mat{R}
  + \Delta\mat{R}\,(-h\mat{B}_n) = h\,\mat{F}_n,
\end{equation}
where $\mat{F}_n = \mat{\alpha}\mat{R}_n + \mat{R}_n\mat{\alpha}
+ \mat{R}_n\mat{\beta}\mat{R}_n + \mat{\beta}$ is the Riccati
right-hand side, and
$\mat{A}_n=\mat{\alpha}+\mat{R}_n\mat{\beta}$,
$\mat{B}_n=\mat{\alpha}+\mat{\beta}\mat{R}_n$.
Equation~\eqref{eq:sylvester_step} is a standard Sylvester equation
$\mat{C}\mat{X}+\mat{X}\mat{D}=\mat{E}$, solvable in $O(N^3)$ via
the Bartels--Stewart algorithm~\cite{bartels1972} (real Schur
decomposition of $\mat{C}$ and $\mat{D}$, followed by
back-substitution).

The companion $\mat{T}$ equation~\eqref{eq:T_ode} has a simpler
implicit step:
$\mat{T}_{n+1}=\mat{T}_n(\mat{I}_N-h\mat{B}_n)^{-1}$,
a single $N\!\times\!N$ linear solve.

\textbf{Cost per step}: one Sylvester solve + one linear solve, both
$O(N^3)$.  No $N^2\!\times\!N^2$ matrix is ever formed.

\subsubsection{Unconditional stability}

The Rosenbrock--Euler method with $\gamma=1$ is L-stable: its
amplification factor $R(z)=1/(1-z)$ satisfies $|R(z)|\to 0$ as
$|z|\to\infty$ in the left half-plane.  For the Riccati ODE, the
fast ballistic modes (eigenvalue $\sim\!-2\kmax$) are immediately
damped, regardless of step size $h$.

\Cref{tab:sylvester} demonstrates this: the method is stable even at
$h\cdot\kmax = 181$, where RK4 and the exponential Euler both fail.

\begin{table}[h]\centering
\caption{Rosenbrock--Euler (Sylvester) on the diffusion sub-domain
($\tsub=20$, $\kmax=45.2$, NQuad${}=16$).  Compare with RK4
stability boundary at $h\cdot\kmax\approx 1.5$ (\Cref{tab:rk4}).}
\label{tab:sylvester}
\begin{tabular}{cccc}
\toprule
$K$ & $h\cdot\kmax$ & Rel.\ err. & Wall (s) \\
\midrule
  5  & 181   & $5.6\times 10^{-2}$ & 0.005 \\
 10  &  90   & $2.4\times 10^{-2}$ & 0.006 \\
 20  &  45   & $1.1\times 10^{-2}$ & 0.014 \\
 50  &  18   & $4.4\times 10^{-3}$ & 0.029 \\
100  &   9   & $2.2\times 10^{-3}$ & 0.054 \\
200  & 4.5   & $1.1\times 10^{-3}$ & 0.10 \\
500  & 1.8   & $4.7\times 10^{-4}$ & 0.22 \\
\bottomrule
\end{tabular}
\end{table}

The convergence is first-order ($O(h)$), as expected for
Rosenbrock--Euler.  At $K=200$ (wall time 0.10~s), the error is
$1.1\times 10^{-3}$---a $14\times$ wall-time speedup over Magnus
($K=914$, wall time 1.43~s) at the cost of lower accuracy.

\subsubsection{Second-order via Richardson extrapolation}
\label{sec:richardson}

Richardson extrapolation of the first-order Rosenbrock--Euler provides
an interim second-order method: solve the Riccati ODE at step counts
$K$ and $2K$, then form
$\mat{R}_{\text{Rich}}=2\mat{R}_{2K}-\mat{R}_K$.
The companion $\mat{T}$ is computed at the finer resolution ($2K$
steps) but is \emph{not} extrapolated, since it depends on the
$\mat{R}(\sigma)$ trajectory.  Total cost: $3K$ Sylvester solves.

\begin{table}[h]\centering
\caption{Richardson extrapolation (cost $3K$ Sylvester solves) on the
diffusion sub-domain ($\tsub=20$, $\kmax=45.2$, NQuad${}=16$).
Compare with Rosenbrock--Euler (\Cref{tab:sylvester}).}
\label{tab:richardson}
\begin{tabular}{cccc}
\toprule
$K$ & $h\cdot\kmax$ & Rel.\ err. & Wall (s) \\
\midrule
  5  & 181   & $6.8\times 10^{-3}$ & 0.005 \\
 10  &  90   & $1.8\times 10^{-3}$ & 0.010 \\
 20  &  45   & $1.0\times 10^{-3}$ & 0.024 \\
 50  &  18   & $6.5\times 10^{-4}$ & 0.053 \\
100  &   9   & $4.2\times 10^{-4}$ & 0.117 \\
200  & 4.5   & $2.5\times 10^{-4}$ & 0.225 \\
500  & 1.8   & $1.2\times 10^{-4}$ & 0.568 \\
\bottomrule
\end{tabular}
\end{table}

At $K=5$ (only 15 Sylvester solves total), Richardson already achieves
$6.8\times 10^{-3}$ error---an $8\times$ improvement over
Rosenbrock--Euler at the same $K$ ($5.6\times 10^{-2}$).  The
improvement is dramatic at large $h\cdot\kmax$ because the Richardson
cancellation removes the dominant $O(h)$ error in $\mat{R}$.

However, the convergence rate degrades as $K$ increases.  From $K=5$
to $K=10$, the error ratio is~$3.8$ (close to the $4\times$ expected
for $O(h^2)$).  From $K=10$ to $K=20$, it drops to~$1.7$---far below
$O(h^2)$.  The reason is that the \emph{companion $\mat{T}$ equation
remains first-order}: $\mat{T}$ is integrated at $2K$ steps (step size
$h/2$) without extrapolation, contributing an $O(h)$ error to the
reconstructed intensity $\vect{u}^+=\mat{R}\vect{b}_{\text{neg}}+
\mat{T}\vect{b}_{\text{pos}}$.  For small $K$, the $O(h^2)$ error in
$\mat{R}$ dominates; for larger $K$, the $O(h)$ error in $\mat{T}$
takes over.

\begin{remark}[Practical comparison at equal accuracy]\label{rem:rich_practical}
For an accuracy target of $\sim\!10^{-3}$: Rosenbrock--Euler requires
$K=200$ (wall time 0.08~s), while Richardson achieves the same accuracy
at $K=20$ (wall time 0.02~s)---a $4\times$ speedup.  The $\mat{T}$
bottleneck prevents further gains beyond this point; a proper two-stage
Rosenbrock method (\Cref{sec:open}, item~6) would give $O(h^2)$ for
both $\mat{R}$ and $\mat{T}$ simultaneously.
\end{remark}

\subsubsection{Comparison of exponential Euler}

An exponential integrator was also tested: the exponential Euler
method with midpoint $\varphi_1$ approximation:
\begin{equation}
  \mat{R}_{n+1} = \mat{R}_n
  + h\,e^{(h/2)\mat{A}_n}\,\mat{F}_n\,e^{(h/2)\mat{B}_n},
\end{equation}
which costs two $N\!\times\!N$ matrix exponentials per step.
This method is A-stable (no blowup) but \emph{not} L-stable:
the fast modes are not actively damped, leading to poor convergence
for $h\cdot\kmax > 10$ (20\% error at $K=50$ vs 0.4\% for
Rosenbrock--Euler).  The L-stability of the Rosenbrock approach is
decisive.

\subsubsection{Scaling with $\omega$}

\Cref{tab:sylvester_omega} shows the Rosenbrock--Euler (Sylvester)
solver with \emph{fixed} $K=200$ across varying $\omega$, compared
with Magnus.

\begin{table}[h]\centering
\caption{Rosenbrock--Euler (Sylvester, $K=200$ fixed) vs Magnus
($K\propto\kmax\tsub$).  Sub-domain $\tsub=20$, $g=0.85$,
NQuad${}=16$.}\label{tab:sylvester_omega}
\begin{tabular}{ccccccc}
\toprule
$\omega$ & $\gamma$ & $K_{\text{Ros}}$ & err$_{\text{Ros}}$
& $t_{\text{Ros}}$ (s) & $K_{\text{M}}$ & $t_{\text{M}}$ (s) \\
\midrule
0.9    &   193 & 200 & $7.4\times 10^{-6}$ & 0.09 & 924 & 1.49 \\
0.99   &   665 & 200 & $1.1\times 10^{-3}$ & 0.10 & 914 & 1.50 \\
0.999  &  2128 & 200 & $2.9\times 10^{-3}$ & 0.10 & 913 & 1.55 \\
0.9999 &  6736 & 200 & $3.2\times 10^{-3}$ & 0.10 & 913 & 1.50 \\
\bottomrule
\end{tabular}
\end{table}

\textbf{Key results}:
\begin{enumerate}[nosep]
\item The Rosenbrock--Euler step count is \emph{independent} of
  $\kmax$ (hence independent of NQuad).  Magnus requires
  $K\propto\kmax\tsub$.
\item Wall time: ${\sim}15\times$ faster than Magnus at NQuad${}=16$.
  The advantage grows with NQuad since $K_{\text{M}}\propto N^2$ while
  $K_{\text{Ros}}$ is fixed.
\item Accuracy degrades as $\omega\to 1$ because $\kmin\to 0$: the
  slow diffusion dynamics require more steps for accurate resolution.
  At $\omega=0.9$ ($\kmin=0.24$), the error is $7\times 10^{-6}$;
  at $\omega=0.9999$ ($\kmin=0.007$), it degrades to $3\times 10^{-3}$.
  Increasing $K$ proportionally to $1/\kmin$ restores accuracy.
\item The method is first-order.  Richardson extrapolation
  (\Cref{sec:richardson}) raises $\mat{R}$ to second order at cost
  $3K$, providing a ${\sim}4\times$ speedup for $10^{-3}$ accuracy.
  However, the companion $\mat{T}$ equation remains first-order,
  limiting overall convergence.  A proper two-stage Rosenbrock method
  giving $O(h^2)$ for both $\mat{R}$ and $\mat{T}$ would overcome
  this bottleneck.
\end{enumerate}

\subsection{Summary of solver comparison}

\Cref{tab:solver_summary} collects the key properties of all four
solvers investigated.

\begin{table}[h]\centering
\caption{Solver comparison for the diffusion sub-domain
($\tsub=20$, $\omega=0.99$, NQuad${}=16$, $N=8$).  ``Rich-Sylv'' is
Richardson extrapolation of Rosenbrock--Euler (\Cref{sec:richardson}).}
\label{tab:solver_summary}
\begin{tabular}{lccccc}
\toprule
& \textbf{Magnus} & \textbf{RK4} & \textbf{Radau}
& \textbf{Ros-Sylv} & \textbf{Rich-Sylv} \\
\midrule
Stability & $h\kmax\lesssim 1$
  & $h\kmax\lesssim 1.5$ & A-stable & L-stable & L-stable \\
Per-step cost & $O(N^3)$ expm & $O(N^3)$ & $O(N^6)$ Jac
  & $O(N^3)$ Sylv & $O(N^3)$ Sylv \\
Steps $K$ & 914 & $\geq 600$ & $\sim 80$ & 200 & 20 \\
Effective evals & 914 & 2400+ & 4820 & 200 & 60 (=$3K$) \\
Rel.\ error & $5\times 10^{-9}$ & $3\times 10^{-3}$ & $5\times 10^{-9}$
  & $1\times 10^{-3}$ & $1\times 10^{-3}$ \\
Wall time (s) & 1.43 & 0.6+ & 0.99 & 0.10 & 0.02 \\
Order ($\mat{R}$/$\mat{T}$) & 2/2 & 4/4 & $\leq$5/$\leq$5 & 1/1 & 2/1 \\
NQuad scaling & $K\propto N^2$ & $K\propto N^2$ & $K$ const
  & $K$ const & $K$ const \\
\bottomrule
\end{tabular}
\end{table}

The Rosenbrock--Sylvester family offers the best wall-time performance
and is the only approach whose step count is truly independent of the
ballistic eigenvalue $\kmax$.  Richardson extrapolation
(\Cref{sec:richardson}) raises $\mat{R}$ to second order at modest
additional cost ($3K$ total Sylvester solves), achieving the same
$10^{-3}$ accuracy as first-order Rosenbrock--Euler but at $K=20$
instead of $K=200$---a $\sim\!4\times$ wall-time improvement.  The
remaining bottleneck is the first-order companion $\mat{T}$ equation:
a proper two-stage Rosenbrock method giving $O(h^2)$ for both
$\mat{R}$ and $\mat{T}$ simultaneously would remove this limitation and
make the Rosenbrock--Sylvester family the clear choice for
diffusion-domain computation.
The Radau solver achieves the highest accuracy per step but pays for
this with the $N^2\!\times\!N^2$ Jacobian cost.


%======================================================================
\section{Open Questions}\label{sec:open}
%======================================================================

\begin{enumerate}
\item \textbf{Automatic detection of diffusion-domain boundaries.}
  The current implementation requires manual specification of
  $\tau_1$ and $\tau_2$.  A criterion based on the beam
  attenuation $e^{-\tau/\mu_0}$ and the local eigenvalue gap
  $\gamma(\tau)$ could be developed.

\item \textbf{$\tau$-varying properties in the implicit Riccati
  solver.}  The current benchmark uses a homogeneous sub-domain.
  For $\tau$-varying $\mat{\alpha}(\tau)$ and $\mat{\beta}(\tau)$,
  the Jacobian~\eqref{eq:jac_riccati} becomes $\sigma$-dependent,
  requiring re-evaluation at each step.

\item \textbf{Interaction with 4th-order Magnus.}
  Report~I~(\S8) developed a 4th-order Magnus extension for smooth
  profiles.  A natural strategy is 4th-order Magnus in the
  non-diffusive domains (where the step size is set by profile
  smoothness) and implicit Riccati in the diffusion domain (where
  stiffness dominates).

\item \textbf{Larger NQuad.}
  For NQuad${}=32$ ($N=16$), the Rosenbrock--Sylvester method
  (\Cref{sec:sylvester}) avoids the $N^2\!\times\!N^2$ Jacobian
  entirely, using $O(N^3)$ Sylvester solves.  However, the constant
  in the $O(N^3)$ cost (real Schur decomposition at each step)
  must be benchmarked against the Magnus $\texttt{expm}$ cost at
  the same $N$.

\item \textbf{$\mat{R}_\downarrow$ and bidirectional Riccati.}
  The current formulation computes only $\mat{R}_\uparrow$.  For a
  homogeneous slab, $\mat{R}_\uparrow=\mat{R}_\downarrow$ by
  symmetry.  For a $\tau$-varying atmosphere, a separate backward
  Riccati integration is needed to obtain $\mat{R}_\downarrow$.

\item \textbf{Second-order Rosenbrock--Sylvester.}
  Richardson extrapolation (\Cref{sec:richardson}) achieves $O(h^2)$
  for $\mat{R}$ but the companion $\mat{T}$ equation remains $O(h)$,
  limiting the overall convergence (\Cref{rem:rich_practical}).  The
  path forward is a proper two-stage L-stable Rosenbrock method (e.g.,
  ROS2~\cite{hairer1996}, Ch.~IV~\S7) that gives $O(h^2)$ for both
  $\mat{R}$ and $\mat{T}$ simultaneously.  Each stage solves one
  Sylvester equation for $\mat{R}$ and one $N\!\times\!N$ linear
  system for $\mat{T}$, so the per-step cost is $2\times$ the
  Rosenbrock--Euler cost.  With the $O(h^2)$ convergence, the step
  count for a given accuracy target would decrease roughly as
  $K\propto 1/\sqrt{\varepsilon}$ (vs $K\propto 1/\varepsilon$ for
  first order), making it the definitive solver for the diffusion
  domain.  The two stages also provide a natural embedded error
  estimate for adaptive step control.

\item \textbf{Coupled $\mat{R}$--$\mat{T}$ integration.}
  The Rosenbrock--Sylvester approach naturally integrates $\mat{R}$
  and $\mat{T}$ in a single pass (the $\mat{T}$ step is a single
  $N\!\times\!N$ solve reusing $\mat{B}_n$ from the Sylvester
  step).  This eliminates the Radau approach's overhead of a
  separate $\mat{T}$ integration pass.

\item \textbf{Symmetry exploitation.}
  For azimuthal mode $m=0$ with Henyey--Greenstein phase functions,
  $\mat{R}$ is symmetric.  Exploiting this reduces the Sylvester
  equation to a \emph{symmetric} Sylvester equation (Lyapunov
  equation), solvable in $O(N^3)$ with a smaller constant.
\end{enumerate}


%======================================================================
\begin{thebibliography}{99}
%======================================================================

\bibitem{report1}
Technical Report~I: ``The Star-Product Magnus Integrator for Radiative
Transfer with Continuously $\tau$-Varying Optical Properties,''
March~2026.

\bibitem{bellman1960}
R.~Bellman, R.~Kalaba, and M.~Prestrud, \emph{Invariant Imbedding and
Radiative Transfer in Slabs of Finite Thickness}, American Elsevier,
1960.

\bibitem{bittanti1991}
S.~Bittanti, A.J.~Laub, and J.C.~Willems (eds.), \emph{The Riccati
Equation}, Springer, 1991.

\bibitem{lancaster1995}
P.~Lancaster and L.~Rodman, \emph{Algebraic Riccati Equations},
Oxford University Press, 1995.

\bibitem{hairer1996}
E.~Hairer and G.~Wanner, \emph{Solving Ordinary Differential
Equations~II: Stiff and Differential-Algebraic Problems}, 2nd~ed.,
Springer, 1996.

\bibitem{bartels1972}
R.H.~Bartels and G.W.~Stewart, ``Solution of the matrix equation
$AX + XB = C$,'' \emph{Comm.\ ACM}, 15(9):820--826, 1972.

\bibitem{efremenko2017}
D.S.~Efremenko, V.~Molina~Garc\'\i a, S.~Gimeno~Garc\'\i a, and
A.~Doicu, ``A review of the matrix-exponential formalism in radiative
transfer,'' \emph{J.\ Quant.\ Spectrosc.\ Radiat.\ Transfer},
196:17--45, 2017.

\bibitem{chandrasekhar1960}
S.~Chandrasekhar, \emph{Radiative Transfer}, Dover, 1960.

\bibitem{stamnes1988}
K.~Stamnes, S.-C.~Tsay, W.~Wiscombe, and K.~Jayaweera, ``Numerically
stable algorithm for discrete-ordinate-method radiative transfer in
multiple scattering and emitting layered media,'' \emph{Appl.\ Opt.},
27(12):2502--2509, 1988.

\bibitem{redheffer1962}
R.~Redheffer, ``On the relation of transmission-line theory to
scattering and transfer,'' \emph{J.\ Math.\ Phys.}, 41(1):1--41, 1962.

\bibitem{grant1969}
I.P.~Grant and G.E.~Hunt, ``Discrete space theory of radiative
transfer.~I.~Fundamentals,'' \emph{Proc.\ R.\ Soc.\ London~A},
313(1513):183--197, 1969.

\bibitem{HP2024doc}
H.~Patel, \emph{Comprehensive Documentation for Pythonic-DISORT},
2024.  Available at
\url{https://pythonic-disort.readthedocs.io}.

\end{thebibliography}

\end{document}
